\documentclass[11pt]{article}

\usepackage[margin=1in]{geometry}
\usepackage{amsmath, amssymb, amsthm}
\usepackage{booktabs}
\usepackage{hyperref}
\usepackage{enumitem}
\usepackage{xcolor}

\newcommand{\kQ}{\kappa_{Q}}
\newcommand{\E}{\mathbb{E}}
\newcommand{\Var}{\operatorname{Var}}
\newcommand{\Cov}{\operatorname{Cov}}
\newcommand{\Rb}{\mathbb{R}}
\newcommand{\Nrm}{\mathcal{N}}
\newcommand{\defA}{\textsf{A}}
\newcommand{\defB}{\textsf{B}}

\title{\textbf{An Error-Decomposition Diagnostic for the\\
Drift-vs-Diffusion Honing Choice}\\[4pt]
\large A design memo: theory context, implemented method, and a
pre-registered decision rule}
\author{MITACS dynamics project --- design memo (not a result artifact)}
\date{2026-05-18}

\begin{document}
\maketitle

\begin{abstract}
We face a fork in honing the predictive estimator. Two deficiencies have
been independently established: \textbf{Deficiency \defA} --- the
conditional \emph{mean} is mis-centred by demand phase (rising vs.\
falling) --- the one candidate deficiency that survived a prior
bootstrap-strength re-test (a separate investigation that retracted
several earlier under-tested claims and established that no
prediction-error cross-validation can recover a localization
bandwidth here, so the drift is effectively a single global linear
map); and \textbf{Deficiency \defB} --- the conditional \emph{law}
is modelled as Gaussian but the true scalar innovation is strongly
heavy-tailed (production-validated excess kurtosis ${\approx}25\text{--}33$).
They live on different axes of the same object and one does not fix the
other. This memo (i) places both deficiencies inside the
Resolvent\_Framework programme's kernel picture, (ii) states precisely
what the implemented estimator does and where each deficiency enters it,
and (iii) specifies a dev-only, variance-checked decomposition diagnostic
with a decision rule \emph{fixed before any numbers are seen}, so the
\defA-vs-\defB\ choice is made on evidence rather than intuition. The
diagnostic is not yet run; this document is its design.
\end{abstract}

\section{Why this memo exists}

The honing thread began from a concrete symptom: on the 2021--22
inspection benchmark the forecast error is not noise --- it is
systematically structured by the daily demand cycle. The investigation
into \emph{why} repeatedly elevated under-tested observations to
``findings'' that later inverted under stricter testing (the prior
investigation referred to above retracted four such over-reaches in
sequence, each from reading a point estimate without a variance or
replication check). The discipline that
emerged, and that this memo obeys, is: \emph{no diagnostic informs a
build decision unless it carries a variance/replication check and a
decision rule pre-registered before the numbers are seen.}

Two candidate deficiencies now survive at different evidentiary
strengths. Choosing which to engineer (or whether to engineer at all) is
not an intuition call. This memo designs the test that decides it.

\section{Theory context: the kernel has a location and a shape}
\label{sec:theory}

\subsection{The programme's kernel}

The Resolvent\_Framework programme (Paper~II) shows that, given a
probability measure on the observable $\sigma$-algebra, the dynamics are
\emph{disclosed} by the measure rather than constructed by the observer.
The Rokhlin disintegration of two observations $Q,F$ yields a Markov
kernel
\[
  \kQ(q,\cdot) \;=\; \mathbb{P}\!\left(F \in \cdot \,\middle|\, Q = q\right),
\]
\emph{forced by the measure, not chosen}. Temporal indexing forces
Chapman--Kolmogorov, yielding a semigroup $\{\Pi_t\}$ and the operator
layer $K_t$. These three layers ($\kQ$, $\{\Pi_t\}$, $K_t$) are kept
notationally distinct. Reconstruction (Paper~II, Part~2) states that if
the observable algebra $\mathcal{O}_h = \sigma(\{h\circ T^n\})$ generates
the full $\sigma$-algebra mod $\mu$, the delay map $\Phi_h$ \emph{is} a
measure-theoretic embedding --- the state space is recovered from the
delayed scalar observable alone.

\subsection{What the code is, in that language}

The implemented estimator is a finite-sample, parametric, single-step,
\emph{locally-Gaussian} estimator of $\kQ$ at the delay-embedding layer.
Per anchor state $x_j$ it produces a pair $(C_j,\Sigma_j)$ defining
\begin{equation}
\label{eq:kernel}
  \widehat{\kQ}(x_j,\cdot) \;=\; \Nrm\!\big(\,x_j C_j,\ \Sigma_j\,\big)
  \qquad\text{on the delay-embedding state space.}
\end{equation}
Three honest qualifiers attach to this correspondence and must not be
dropped: \textbf{(Q1)} register --- the programme \emph{discloses} $\kQ$
from the measure; the code \emph{constructs} a Gaussian proxy and fits it
(``estimator of,'' not ``instance of''); \textbf{(Q2)} Gaussianity is
\emph{imposed}, exact only where the local conditional law is actually
Gaussian; \textbf{(Q3)} single-step only --- $\{\Pi_t\}$ is neither
constructed nor verified, so this is the slice $\Pi_\Delta$, not the
semigroup.

\subsection{The fork, stated in the kernel}
\label{sec:fork}

Equation~\eqref{eq:kernel} makes the fork precise. A conditional law has,
at minimum, a \emph{location} and a \emph{shape}. The estimator
factorises these into two distinct objects:
\begin{itemize}[leftmargin=2em]
  \item the \textbf{drift} $C_j$ sets the conditional \emph{location}
        $\E[x' \mid x_j] = x_j C_j$;
  \item the \textbf{diffusion} $\Sigma_j$ (with the imposed Gaussian
        form) sets the conditional \emph{shape/spread}.
\end{itemize}
The two established deficiencies are deficiencies of these two distinct
factors:

\begin{center}
\begin{tabular}{@{}p{0.10\linewidth}p{0.40\linewidth}p{0.40\linewidth}@{}}
\toprule
 & \textbf{Deficiency \defA} & \textbf{Deficiency \defB}\\
\midrule
Factor & Drift $C_j$ (location) & Diffusion $\Sigma_j$ (shape) \\[3pt]
Claim  & The location is \emph{mis-placed}: a single global-ish
         linear map cannot separate rising-through-a-level from
         falling-through-a-level, so it centres between them.
       & The shape is \emph{mis-modelled}: a Gaussian is imposed on a
         scalar innovation that is strongly heavy-tailed. \\[3pt]
Programme reading
       & The conditional-mean estimate of $\kQ$ is biased on a
         phase coordinate. Same target $\kQ$, location wrong.
       & Directly qualifier \textbf{Q2}: $\Sigma_j$ is a Gaussian
         second-moment \emph{proxy} of a non-Gaussian law. Same target
         $\kQ$, shape wrong. \\[3pt]
Evidence
       & \textbf{The phase-structure finding}: the signed forecast
         error correlates with the demand derivative,
         $\mathrm{corr}(\mathrm{d}z/\mathrm{d}t,\ z_{\mathrm{err}})
         = -0.246$, bootstrap 95\% CI $[-0.283,-0.207]$
         ($n=3167$) --- survived the same bootstrap-strength
         re-test that retracted the other honing leads.
       & Production re-baseline: intra-day scalar-innovation excess
         kurtosis ${\approx}24\text{--}33$, tail-ratio
         ${\approx}2.4\text{--}3.8$ (Gaussian $=0$ and $1$; clean
         VAR(1) reference ${\approx}0$ and ${\approx}1$). \\
\bottomrule
\end{tabular}
\end{center}

\noindent\textbf{Crucially, neither fixes the other.} A heavy-tailed
innovation law leaves $x_j C_j$ exactly where it was --- it does not
re-centre a mis-placed conditional mean. A phase-aware drift sharpens the
location but says nothing about whether the residual law around it is
Gaussian. \defA\ and \defB\ are orthogonal in the factorisation of
\eqref{eq:kernel}. Both are ``same-$\kQ$'' estimator-class improvements
(neither changes $\Phi_h$ or $\mathcal{O}_h$, so the reconstruction
object is untouched); they differ only in \emph{which factor} of the
Gaussian proxy is made more honest. That is exactly why the choice
between them needs evidence about \emph{which factor's error dominates
the forecast loss we care about}.

\section{The implemented method (what we are decomposing the error of)}
\label{sec:method}

The diagnostic must read the \emph{production} estimator, not a
re-implementation (the project's hardest-won discipline: a diagnostic
that re-implements production logic is a hypothesis, not a finding, until
confirmed through the production path). The relevant production facts:

\begin{enumerate}[leftmargin=2em]
  \item \textbf{Embedding.} A single-variable delay embedding of the
        de-seasonalised $z$-score, $d$ lags per day-type (the elbow-rule
        dimensions). $\Sigma_j$ is therefore \emph{rank-1 by
        construction}: coordinates $2..d$ of the one-step image are
        deterministic shifts of the input, so the only stochastic content
        is the scalar coordinate-0 innovation
        $r_0 = z_{t+1} - (xC)_0$. Non-Gaussianity (\defB) is correctly a
        \emph{univariate} question on $r_0$.
  \item \textbf{Drift $C_j$.} Weighted least squares with a normalised
        Gaussian kernel; the bandwidth $\theta_j$ is selected per anchor
        by true leave-one-out one-step CV. The prior investigation
        established that this CV objective has no statistically
        separated interior optimum on any system tested, so in effect
        the drift is a
        \emph{single global-ish linear map}. Deficiency \defA\ is a
        statement about \emph{this map's} conditional-mean placement.
  \item \textbf{Diffusion $\Sigma_j$.} The $\mu$-centred \emph{plain}
        covariance of the locally-fitted residuals $Y - XC$ (no residual
        kernel). It is the unbiased local second moment --- and, under
        Gaussianity, the maximum-likelihood shape. Deficiency \defB\
        says this second-moment object is a Gaussian stand-in for a
        heavy-tailed law.
  \item \textbf{Estimand.} The intra-day-conditioned $\kQ$: full-process
        $\Sigma_j$ is numerically unusable (the day-anchor seam wrecks
        conditioning), so transitions whose one-step target lands on the
        day-anchor hour are masked. All numbers here are intra-day.
  \item \textbf{Use.} Production forecasting iterates the one-step
        kernel into a 24-step daily trajectory. The error we ultimately
        care about is the \emph{multi-step} forecast error on the
        2021--22 dev benchmark --- which is what the phase-structure
        finding was measured on, and what this decomposition must
        target.
\end{enumerate}

\section{The diagnostic: decomposing forecast loss into location vs.\ shape}
\label{sec:diag}

\subsection{Principle}

For each scored dev forecast point we have a predictive law
$\widehat{\kQ}(x,\cdot)=\Nrm(\,\widehat{m},\widehat{s}^{\,2})$ for the
scalar quantity of interest (coordinate-0 $z$, mapped back as needed) and
a realised actual $y$. The forecast loss decomposes, in expectation, into
a part driven by \emph{where the law is centred} and a part driven by
\emph{the law's shape/spread}. We make this operational two ways and
require them to agree before drawing a conclusion (cross-method
replication is itself part of the discipline).

\paragraph{Obtaining $\widehat{s}^{\,2}$ at every horizon (resolved).}
The production estimator emits only the \emph{one-step} local diffusion
$\Sigma$; the benchmark forecast is iterated to $24$ horizons, so the
predictive \emph{law} at horizon $k$ is not $\Nrm(\widehat{m}_k,
\Sigma[0,0])$. We therefore propagate the predictive covariance through
the iteration: with the shift-aware state Jacobian
$J_i = [\,C_i[:,0]\mid\text{shift block}\,]$ and the rank-1 per-step
innovation $\Sigma_{\text{step},i}=\mathrm{diag}(\Sigma_i[0,0],0,\dots)$
(the true iterated-forecast semantics; the rank-1 form is structural,
not assumed), the predictive covariance obeys
$P_{i+1}=J_i^{\!\top}P_iJ_i+\Sigma_{\text{step},i}$, $P_0=0$, and
$\widehat{s}^{\,2}_k=P_k[0,0]$. This composition is the programme's
$\{\Pi_t\}$ layer, which qualifier \textbf{Q3} flags as not previously
constructed; it is therefore \emph{gate-validated against the VAR(1)
closed form $\Sigma_k=\sum_{i<k}(A^{\!\top})^iQA^i$ to machine precision
(including the production rank-1 shift-Jacobian shape) before it is
allowed to inform this diagnostic} --- the same validate-before-use bar
the estimator itself meets. Linearisation around the realised path is
the standard EKF-style scheme; for a linear local model it is exact
within the iteration, the only unpropagated term being $C_i$'s own
dependence on the local state, which is exactly what the production
point forecast also does.

\paragraph{The embedding-dimension-change boundary (a real fork).}
The composition $P_{i+1}=J_i^{\!\top}P_iJ_i+\Sigma_{\text{step},i}$ is
only \emph{defined} when the state dimension is constant across the
step. It is not: the per-day-type embedding dimensions differ
($\text{weekday}=2$, $\text{saturday}=4$, $\text{sunday}=3$) and the
day-type rolls at the $07{:}00$ anchor, so \textbf{${\approx}42\%$ of
delivery days cross a dimension change within their $24$-hour forecast}
(the \emph{point} forecast is unaffected --- it rebuilds the lag state
each step; only the covariance composition $J^{\!\top}PJ$ breaks across
$d_i\neq d_{i+1}$). There is no closed-form ground truth for the
predictive covariance across an embedding-induced dimension change, so
this is genuinely a choice, not a derivation. Four handlers are
considered; only the first is a validated, trustworthy verdict, the
other three are \emph{exploratory robustness probes} (unvalidated
approximations, inspection-only) whose sole purpose is to test whether
the conclusion depends on this choice:

\begin{center}
\begin{tabular}{@{}p{0.13\linewidth}p{0.50\linewidth}p{0.30\linewidth}@{}}
\toprule
\textbf{Scope} & \textbf{Boundary handling} & \textbf{Status}\\
\midrule
\texttt{clean}
 & Restrict to delivery days whose entire $24$h forecast stays at one
   dimension; $P_{i+1}=J^{\!\top}P_iJ+\Sigma_{\text{step}}$ is exactly
   defined throughout.
 & \textbf{The verdict.} No approximation, no statable bias. Pays an
   exclusion of the ${\approx}42\%$ rollover days (caveat below). \\[10pt]
\texttt{reset}
 & On a dimension change set $P\!\to\!0$ and re-accumulate.
 & Exploratory. \emph{Known signed bias}: understates
   $\widehat{s}^{\,2}_k$ after any seam (discards accumulated forecast
   uncertainty) --- inflates apparent calibration on the shape axis. \\[10pt]
\texttt{project}
 & On a dimension change map $P$ across by truncation (dim drop) /
   zero-pad (dim rise) of the leading block.
 & Exploratory. An \emph{assumption} about cross-dimension state
   covariance with no ground truth; cannot be gate-validated as the
   in-dimension recursion was. \\[10pt]
\texttt{h1}
 & Ignore propagation; $\widehat{s}^{\,2}=\Sigma_0[0,0]$ (the exact
   one-step law).
 & Exploratory. No boundary issue at all, but tests only horizon~$1$,
   not the multi-step error the phase-structure finding was measured
   across. \\
\bottomrule
\end{tabular}
\end{center}

\noindent\textbf{Exclusion caveat (load-bearing, travels with the
verdict).} The \texttt{clean} scope drops every day-type-rollover day.
If Deficiency \defA\ (phase mis-centring) \emph{concentrates at
day-type seams} --- plausible, since seams are where the asymmetric
ramp is most violent --- the \texttt{clean} verdict \emph{under-counts}
\defA. The conclusion is therefore conditional on the single-dimension
subpopulation. The three exploratory scopes exist precisely to probe
this: if every handler yields the same A-vs-B verdict, the exclusion is
materially de-risked; if they diverge, the exclusion is load-bearing
and the divergence is itself the finding. Either way, only
\texttt{clean} is reported as a validated verdict --- the probes are
sensitivity context, never elevated to a result (the project's
inspection-only discipline; per the recording decision they are recorded
in the verdict memo with this caveat, not stamped through
\texttt{make\_result}).

\paragraph{(M1) Mean/variance decomposition of squared error.}
With signed error $e=\widehat{m}-y$, write the conditional mean-squared
error as
\[
  \E[e^2 \mid \text{stratum}]
   \;=\; \underbrace{\big(\E[e\mid\text{stratum}]\big)^2}_{\text{bias}^2\ \to\ \defA}
   \;+\; \underbrace{\Var[e\mid\text{stratum}]}_{\text{dispersion}\ \to\ \defB}.
\]
Stratify by the phase coordinate $\mathrm{d}z/\mathrm{d}t$ (the same
demand-derivative coordinate the phase-structure finding used; in
quintiles) and by hour-of-day. A large,
phase-structured $\mathrm{bias}^2$ fraction implicates the
\emph{location} factor (\defA); a dispersion term whose \emph{magnitude}
is dominated by rare large $|e|$ implicates the \emph{shape} factor
(\defB). The bias term is the \defA\ signature; the bias-removed loss is
the ceiling that only a shape change (\defB) could further reduce.

\paragraph{(M2) Counterfactual log-loss / CRPS under each fix.}
Evaluate the predictive law's proper score (Gaussian log-loss and CRPS)
on the dev set under three predictive laws built \emph{only from
production outputs}:
\begin{enumerate}[label=(\alph*),leftmargin=2.2em]
  \item \textbf{baseline}: $\Nrm(\widehat{m},\widehat{s}^{\,2})$ as
        produced;
  \item \textbf{\defA-counterfactual}: subtract the dev-estimated
        phase-conditional bias from $\widehat{m}$ (an \emph{oracle}
        upper bound on what perfect re-centring could buy --- not a
        proposed fix, a bound), keep $\widehat{s}^{\,2}$;
  \item \textbf{\defB-counterfactual}: keep $\widehat{m}$, replace the
        Gaussian shape by a heavy-tailed fit (Student-$t$ with dev-fitted
        $\nu$) matched to $\widehat{s}^{\,2}$.
\end{enumerate}
The score improvement of (b) over (a) bounds the \defA\ lever; of (c)
over (a) bounds the \defB\ lever, \emph{on the loss the production system
is actually graded on}. These are diagnostic counterfactuals, not the
builds themselves --- they answer ``which axis has more recoverable
loss,'' nothing more.

\subsection{Pre-registered decision rule (fixed before any numbers)}
\label{sec:rule}

The statistic is the \textbf{bootstrapped difference in dev forecast
loss} between each counterfactual and baseline, resampling whole forecast
\emph{days} with replacement ($B=1000$; day-level resampling respects the
24-step iteration dependence). Let
$\Delta_\defA$ and $\Delta_\defB$ be the per-bootstrap loss reductions
(positive $=$ improvement) under M2, with 95\% percentile CIs; M1's
$\mathrm{bias}^2$ fraction is reported with a day-bootstrap CI as the
corroborating second method.

\begin{center}
\begin{tabular}{@{}p{0.30\linewidth}p{0.62\linewidth}@{}}
\toprule
\textbf{Outcome} & \textbf{Pre-registered conclusion}\\
\midrule
$\Delta_\defA$ CI lower bound $>0$ \emph{and}
excludes $\Delta_\defB$'s CI (non-overlapping, $\defA$ larger)
 & \defA\ dominates: the \emph{drift/location} lever is the
   higher-leverage honing target. Pursue a phase-aware conditional mean;
   \defB\ remains a documented Q2 caveat. \\[4pt]
$\Delta_\defB$ CI lower bound $>0$ \emph{and}
excludes $\Delta_\defA$'s CI (non-overlapping, $\defB$ larger)
 & \defB\ dominates: the \emph{diffusion/shape} lever is higher-leverage.
   Pursue dropping local-Gaussian (heavy-tailed innovation law); \defA\
   becomes a documented limitation. \\[4pt]
CIs overlap, \emph{or} neither lower bound clears
$1\%$ relative loss reduction
 & \textbf{Neither lever is decisively supported.} Default to the
   standing honest-characterisation option: document \emph{both} \defA\
   and \defB\ as limitations of the univariate local-Gaussian $\kQ$
   estimator and proceed to the registered forward experiment unchanged.
   (This is the project's repeatedly-chosen ``characterise, don't
   over-engineer'' resolution; the prior bandwidth investigation
   ended the same way.) \\
\bottomrule
\end{tabular}
\end{center}

\noindent The relative-loss threshold ($1\%$ of baseline dev loss) and
the non-overlapping-CI requirement are fixed here, before execution. M1
and M2 must \emph{agree in direction}; if they disagree the result is
declared inconclusive (not adjudicated post-hoc) and the disagreement
itself is the finding.

\subsection{What this diagnostic does \emph{not} establish}

Honest scope, mirroring the phase-structure finding's own caveat:
\begin{itemize}[leftmargin=2em]
  \item It bounds the \emph{recoverable loss} on each axis via oracle
        counterfactuals. It does \emph{not} prove that any \emph{specific}
        buildable model (e.g.\ piecewise-linear-in-$\mathrm{sign}(\mathrm{d}z/\mathrm{d}t)$
        drift, or a Student-$t$ innovation) attains that bound. The
        counterfactual is an upper bound on the lever, not a guarantee a
        feasible estimator reaches it.
  \item It is a \emph{dev-set} measurement (2021--22, $\le$ cutoff fit /
        dev-window evaluation). Per the integrity rule, any model change
        it motivates must be re-derived dev-only and locked into a new
        hash-stamped frozen spec before any forward scoring.
  \item It does not re-open the embedding-layer ($\Phi_h$) or
        observable-algebra ($\mathcal{O}_h$) options; by construction
        both counterfactuals are same-$\kQ$ estimator-class moves. The
        scope-expanding option (exogenous covariates) is deliberately
        out of frame here.
\end{itemize}

\section{Results}
\label{sec:results}

\subsection{\texttt{clean} scope --- the validated verdict}

Run on the regenerated 2021--22 dev benchmark with the gate-validated
variance propagation, \texttt{clean} scope (single-embedding-dimension
delivery days: $76$ kept, $56$ day-type-rollover days excluded),
day-level bootstrap $B=1000$:

\begin{center}
\begin{tabular}{@{}lll@{}}
\toprule
\textbf{Statistic} & \textbf{Estimate} & \textbf{95\% CI} \\
\midrule
M1: bias$^2$ fraction of MSE & $0.377$ & $[0.358,\ 0.484]$ \\
M2: \defA-counterfactual (oracle re-centre)
   & $+6.11\%$ rel.\ log-loss & $[+3.42\%,\ +8.72\%]$ \\
M2: \defB-counterfactual (Student-$t$, $\nu\approx4.98$)
   & $+9.04\%$ rel.\ log-loss & $[+5.44\%,\ +12.81\%]$ \\
\multicolumn{2}{@{}l}{M2: fitted $\nu$ / std.-resid.\ excess kurtosis}
   & $4.98$ / $6.14$ \\
\bottomrule
\end{tabular}
\end{center}

\noindent\textbf{Pre-registered verdict: NEITHER dominates.} Both
oracle counterfactuals clear the $1\%$ relative-loss bar and are
statistically separated from zero --- so \emph{both} Deficiency \defA\
and \defB\ are real and non-trivial. But the two CIs \emph{overlap}
($[3.4,8.7]\%$ vs.\ $[5.4,12.8]\%$), so the pre-registered
non-overlapping-CI requirement for declaring a dominant lever is not
met; M1 corroborates ($0.38$ bias-fraction, neither CI clearing the
$0.5$ line --- it points at \emph{both}, decisively at \emph{neither}).
This is \emph{not} ``no deficiency'': it is ``both deficiencies are
real and comparable in magnitude; the evidence does not separate them,''
and $\defB$'s central estimate is the larger of the two. The
pre-registered outcome is therefore to document \emph{both} \defA\ and
\defB\ as quantified limitations of the univariate local-Gaussian
$\kQ$ estimator, build nothing, and proceed to the registered forward
experiment unchanged.

\subsection{Robustness sweep over the dimension-change scopes}

To test whether ``neither dominates'' is an artifact of the
${\approx}42\%$ rollover-day exclusion, the diagnostic was first re-run
under three \emph{exploratory} boundary handlers (\texttt{reset},
\texttt{project}, \texttt{h1}; unvalidated approximations, no
closed-form ground truth) and then, decisively, under a \emph{validated}
all-days path (\texttt{augmented}). The pre-registered reading rule for
the exploratory sweep was binary and fixed before its numbers were
seen: \emph{if any exploratory scope diverges from \texttt{clean}, the
exclusion is load-bearing and only the validated path settles it}. The
exploratory sweep \textbf{failed} that rule (\texttt{h1} diverged) ---
which is precisely why the validated \texttt{augmented} path was built;
its result (top row) is the resolution and is discussed after the
sweep.

\begin{center}
\begin{tabular}{@{}lllll@{}}
\toprule
\textbf{Scope} & \textbf{M1 bias$^2$ (CI)}
 & \textbf{M2 \defA\ rel.\ (CI)} & \textbf{M2 \defB\ rel.\ (CI)}
 & $\nu$ / \textbf{Verdict} \\
\midrule
\texttt{augmented} (\textbf{valid, all days})
 & $0.330\ [.30,.42]$ & $+4.8\%\ [0.7,8.6]$
 & $+2.3\%\ [-0.0,4.8]$ & $6.48$ / \textbf{NEITHER} \\
\texttt{clean} (\emph{valid, single-dim})
 & $0.377\ [.36,.48]$ & $+6.1\%\ [3.4,8.6]$
 & $+9.1\%\ [5.5,12.9]$ & $4.98$ / NEITHER \\
\texttt{reset} \emph{(expl.)}
 & $0.330\ [.29,.43]$ & $+15.1\%\ [8.6,20.7]$
 & $+21.8\%\ [16.7,26.6]$ & $4.51$ / NEITHER \\
\texttt{project} \emph{(expl.)}
 & $0.330\ [.30,.42]$ & $+5.1\%\ [0.9,9.1]$
 & $+2.2\%\ [-0.1,4.5]$ & $6.60$ / NEITHER \\
\texttt{h1} \emph{(expl.)}
 & $0.330\ [.30,.42]$ & $+15.7\%\ [9.2,21.4]$
 & $+28.7\%\ [25.9,31.6]$ & $57.95$ / \textbf{B} \\
\bottomrule
\end{tabular}
\end{center}

\noindent\textbf{The sweep FAILED its pre-registered robustness rule:
the exclusion is load-bearing.} The rule fixed in advance was binary
--- if \emph{any} scope diverges from \texttt{clean}, the
rollover-day exclusion is load-bearing and only \texttt{clean} stands.
\texttt{h1} diverges (``\defB\ dominates''). Per the pre-registered
rule the outcome is therefore \textbf{NOT robust}; the divergence
\emph{is} the finding. (An earlier draft of this subsection re-narrated
the sweep as ``robust where it matters'' by explaining \texttt{h1}
away --- that was a post-hoc rationalisation that failed its own
arithmetic and is retracted in full; it is recorded here as a worked
example of exactly the failure mode the pre-registered rule exists to
block.)

\emph{Why the retracted explanation was wrong.} It claimed
$\nu\approx58$ made the \defB\ counterfactual ``nearly vacuous'' at
\texttt{h1} --- but \texttt{h1}'s \defB\ effect is $+28.7\%$, the
\emph{largest} of any scope; a vacuous counterfactual cannot be the
largest. The true mechanism: \texttt{h1}'s $\widehat{s}^{\,2}_k$
(median ${\approx}0.015$) underestimates the realised residual
variance (${\approx}0.27$) by ${\approx}17\times$, so the baseline
Gaussian log-loss is catastrophic (${\approx}8.2$ vs ${\approx}0.67$
at \texttt{clean}); the variance-matched Student-$t$ merely rescues
that gross overconfidence. \texttt{h1}'s ``\defB\ dominates'' is an
artefact of a degenerate predictive variance, \emph{not} evidence
about Deficiency \defB. (The companion claim that \defA\ has ``almost
no phase bias at \texttt{h1}'' was also false: \defA\ at \texttt{h1} is
$+15.7\%$, larger than \texttt{clean}'s $+6.1\%$.)

\emph{And the three ``agreeing'' scopes do not corroborate.}
\texttt{reset} ($\defB=+21.8\%$) and \texttt{project}
($\defB=+2.2\%$, CI $[-0.1,4.5]$ including zero) disagree on the \defB\
magnitude by ${\approx}10\times$ via opposite internal stories; they
share the ``NEITHER'' label only because it is the catch-all bin
(returned on CI-overlap \emph{or} sub-$1\%$). Shared-label is not
shared-evidence: the agreement carries ${\approx}$zero probative
weight.

\noindent\textbf{Resolution: the load-bearing exclusion was
subsequently closed by the validated all-days path.} The conclusion
above (only \texttt{clean} validated; the exclusion load-bearing; the
only honest route is a propagation that genuinely handles the
dimension-change boundary) stood as written --- and that propagation
was then built. A fixed-$d_{\max}$ state augmentation
(\texttt{augmented}) carries the lower-dimensional day-types' unused
lag coordinates as deterministic, zero-innovation components, so the
covariance recursion is exactly defined across a dimension change. It
was validated to the project's standard before use: machine-precision
agreement with the VAR(1) closed form on the active block, and
\emph{bit-identical} ($\max|\Delta|=2.8\times10^{-15}$) to the
existing single-dimension recursion on single-dimension days. (A
cheaper scalar-recursion candidate was \emph{falsified first} by a
deterministic known-answer check --- the propagated covariance is not
rank-one, so a scalar channel misses cross-term feedback and diverges
by the third step --- and was discarded rather than shipped.)

Re-running the diagnostic over \emph{all $132$ delivery days} through
this validated path (the \texttt{augmented} row above) yields
\textbf{NEITHER dominates}, in agreement with the single-dimension
\texttt{clean} verdict. A discriminating per-row check confirms
\texttt{augmented} is genuinely distinct from the unvalidated
\texttt{project} probe (identical to machine precision on
single-dimension days as the gate guarantees, but differing by up to
$7\times10^{-2}$ on dimension-change days --- the close aggregate is
coincidental averaging, not the same arithmetic). The
${\approx}42\%$ exclusion is therefore \textbf{resolved, not merely
caveated}: the validated full-population result and the validated
sub-population result agree. The earlier worry that Deficiency \defA\
might concentrate at day-type seams is not borne out --- including the
seam days slightly \emph{lowers} the apparent \defA\ effect
($+4.8\%$ vs.\ $+6.1\%$), it does not raise it.

\textbf{The two ``NEITHER''s are different evidentiary states --- the
full-population one is stronger.} On \texttt{clean} both effects'
CI lower bounds clear the pre-registered $1\%$ materiality bar
($3.4\%$, $5.5\%$): ``NEITHER'' there means \emph{both real but
statistically tied}. On the validated \texttt{augmented} full
population \emph{neither} lower bound clears the bar (\defA\ $0.7\%$,
\defB\ $-0.0\%$, the \defB\ interval including zero): ``NEITHER'' there
is the \emph{catch-all} --- neither effect is materially established at
all. Same label, stronger conclusion: the evidence does not establish
either deficiency as actionable on the complete, correctly-propagated
data. Relatedly, the heavy-tail signal attenuated ${\approx}4\times$
(Student-$t$ $\nu\approx5.0\to6.5$; implied standardised-residual
excess kurtosis ${\approx}6.1\to2.4$) once predictive variance was
\emph{propagated} across the seam rather than sidestepped by
exclusion --- indicating part of the previously-reported strong
heavy tail (one-step re-baseline excess kurtosis ${\approx}24$--$33$,
measured without multi-step variance propagation) was a propagation
artefact, not solely an innovation-law property. The substantive
outcome is unchanged and strengthened on the full population:
document both \defA\ and \defB\ as quantified limitations of the
univariate local-Gaussian $\kQ$ estimator, build nothing, proceed to
the
registered forward experiment unchanged. \texttt{augmented} and
\texttt{clean} are both validated; \texttt{reset}/\texttt{project}/%
\texttt{h1} remain superseded inspection-only probes (\texttt{h1}'s
``\defB\ dominates'' is the degenerate-predictive-variance artefact
diagnosed above). All rows are inspection-only, recorded in the
verdict memo, never stamped through \texttt{make\_result}.

\section{Summary}

The fork is not ``which feature to add.'' In the programme's kernel
picture it is: \emph{is the Gaussian proxy of $\kQ$ failing more at its
location ($C_j$, Deficiency \defA) or at its shape ($\Sigma_j$,
Deficiency \defB)?} Both are honest same-target estimator improvements;
they are orthogonal and one cannot substitute for the other. This memo
specifies a dev-only, production-path, day-bootstrapped, cross-validated
(M1$\,\wedge\,$M2) decomposition with a decision rule fixed in
advance --- including an explicit, pre-committed ``neither dominates
$\Rightarrow$ document both and do not engineer'' branch, because that is
a real and defensible outcome, not a failure of the test. Nothing is
built until this returns a verdict under its own pre-registered standard.

\end{document}
